\documentclass[11pt]{article}
\usepackage[T1]{fontenc}
\usepackage{amsmath,amssymb,amsthm}
\usepackage[margin=1in]{geometry}

\newtheorem{theorem}{Theorem}
\newtheorem{lemma}{Lemma}
\newtheorem{proposition}{Proposition}
\newtheorem{remark}{Remark}

\title{Laboratory V81: Deficiency Conservation and the Minimum Union Boundary}
\author{NC0\_k-Avoid Laboratory}
\date{August 2026}

\begin{document}
\maketitle

Let $M$ be the output gates, let $X=N(M)$ be the active input variables, and
write $\sigma=|M|-|X|$.  For $S\subseteq M$, define
\[
 \delta(S)=|S|-|N(S)|
 \quad\text{and}\quad
 \lambda_C(S)=|N(S)\cap N(M\setminus S)|.
\]

\begin{theorem}[Deficiency conservation]
For every $S\subseteq M$,
\[
 \delta(S)+\delta(M\setminus S)=\sigma-\lambda_C(S).
\]
\end{theorem}

\begin{proof}
Inclusion--exclusion gives
\[
 |N(S)|+|N(M\setminus S)|=|X|+\lambda_C(S).
\]
Subtracting this equality from $|S|+|M\setminus S|=|M|$ proves the claim.
\end{proof}

\begin{lemma}[Balanced edge]
Every unrooted subcubic tree with $m\geq2$ labelled leaves has an edge whose
two leaf sides have sizes between $\lceil m/3\rceil$ and
$\lfloor2m/3\rfloor$.
\end{lemma}

\begin{proof}
Orient an edge toward a side containing more than $2m/3$ leaves whenever such
a side exists.  At a vertex with no outgoing incident edge, every incident
component has at most $2m/3$ leaves.  At most three components meet there and
their sizes sum to $m$, so one has at least $m/3$ leaves.  Its incident edge
has the required balance.
\end{proof}

\begin{theorem}[Width--deficiency tradeoff]
Given a support branch decomposition of width at most $w$, one can find in
polynomial time a balanced gate set $S$ satisfying
\[
 \delta(S)\geq \left\lceil\frac{\sigma-w}{2}\right\rceil.
\]
When the right-hand side is nonpositive the assertion is only a numerical
bound, not a Hall-deficiency claim.
\end{theorem}

\begin{proof}
Scan the decomposition edges and choose the balanced edge from the lemma.
Its displayed cut has $\lambda_C(S)\leq w$.  Conservation yields
\[
 \delta(S)+\delta(M\setminus S)\geq\sigma-w.
\]
One of the two integral deficiencies is at least the ceiling of half this
quantity; output that side.
\end{proof}

For $0\leq p\leq |M|$, define
\[
 U(p)=\min\{|N(S)|:S\subseteq M,\ |S|=p\}.
\]

\begin{proposition}[Minimum union equivalence]
The minimum neighborhood size of a Hall-deficient set equals
\[
 \min\{U(p):U(p)<p\}.
\]
\end{proposition}

\begin{proof}
Every deficient set of cardinality $p$ witnesses $U(p)<p$.  Conversely, a
set attaining any $U(p)<p$ is Hall deficient.
\end{proof}

\begin{remark}[Lagrangian limitation]
For fixed rational $\alpha\geq0$, the function
$|N(S)|-\alpha|S|$ is submodular and can be minimized as a minimum-weight
closure problem.  Varying $\alpha$ exposes only supported points of the lower
curve $(p,U(p))$.  Exact V81 enumeration shows that the minimum-neighborhood
Hall points of all three V80 families are unsupported; in each case the only
supported deficient point is the full gate set.  Thus the Lagrangian scan is
not a complete algorithm for the constrained problem.
\end{remark}

\begin{remark}[Scope]
The theorem does not derandomize Range Avoidance, solve Minimum $p$-Union,
prove an all-orders lower bound, activate a direct route to P versus NP, or
resolve P versus NP.  Novelty and peer review are not claimed.
\end{remark}

\paragraph{References.}
E.~Chlamtac, M.~Dinitz, C.~Konrad, G.~Kortsarz, and G.~Rabanca,
\emph{The Densest $k$-Subhypergraph Problem}, arXiv:1605.04284.
E.~Chlamtac, M.~Dinitz, and Y.~Makarychev,
\emph{Minimizing the Union: Tight Approximations for Small Set Bipartite
Vertex Expansion}, arXiv:1611.07866.

\end{document}
